
\chapter{Исследовательская часть}

В данном разделе проводится исследование зависимости времени отрисовки одного кадра от разрешения квадратного изображения с длиной стороны N пикселей. 
Замеры времени будут проводиться на одной и той же сцене, позиция камеры охватывает весь город. В качестве значения будет использоваться усреднённое время отрисовки за 50 кадров.

Исследования производились на машине со следующими характеристиками:

\begin{itemize}
	
	\item операционная система --- Arch Linux based on Linux 6.17.8~\cite{arch_wiki};
	
	\item процессор --- AMD Ryzen 7 8750H (16)~\cite{ryzen};
	
	\item оперативная память --- DDR5 6400МГц 24 ГБ.
	
\end{itemize}

Во время замеров ноутбук был подключен к питанию, а также был активирован производительный режим.

\section{Результаты замеров}

При тестировании использовался город квадратной формы. Условные единицы, используемые как мера длины в программе, примерно равны метру.

\input{tables/issled}

\begin{figure}[H]
	\centering
	\includegraphics[width=1\textwidth]{issled}
	\caption{Зависимость времени отрисовки квадратного изображения от размера стороны}
	\label{fig:issled}
\end{figure}

Линия графика показывает квадратичную аппроксимацию по точкам, полученным в результате проведения замеров.

Проведенные исследования показывают, что время отрисовки имеет квадратичную зависимость от стороны квадратного изображения, однако, при городе размером в 1000000 есть увеличенное константное время, обусловленное перебором большого числа полигонов в процессе отображения.

\section*{Вывод}
В результате замеров выяснилось, что при квадратном изображении время отрисовки кадра квадратично зависит от длины стороны. При наличии большого числа полигонов существует константная задержка, обусловленная перебором всех полигонов.